%---------- Inleiding ---------------------------------------------------------

% TODO: Is dit voorstel gebaseerd op een paper van Research Methods die je
% vorig jaar hebt ingediend? Heb je daarbij eventueel samengewerkt met een
% andere student?
% Zo ja, haal dan de tekst hieronder uit commentaar en pas aan.

%\paragraph{Opmerking}

% Dit voorstel is gebaseerd op het onderzoeksvoorstel dat werd geschreven in het
% kader van het vak Research Methods dat ik (vorig/dit) academiejaar heb
% uitgewerkt (met medesturent VOORNAAM NAAM als mede-auteur).
% 

\section{Inleiding}%
\label{sec:inleiding}

Grote organisaties in sectoren zoals financi\"ele dienstverlening, verzekeringen en overheid blijven IBM mainframes inzetten voor bedrijfskritische processen, onder meer door hun hoge beschikbaarheid en betrouwbaarheid voor transactionele workloads \autocite{IBMZArchitecture}. 
Binnen deze context vormt z/OS een kernplatform omdat het ontworpen is voor continue beschikbaarheid, sterke beveiligingsmechanismen en grootschalige verwerking van transacties en batchverwerking \autocite{IBMzOSOverview,IBMRedbookzOS}. 

Het meest gebruikte besturingssysteem op deze platforms, z/OS, is specifiek ontworpen om een stabiele, veilige en continu beschikbare omgeving te bieden voor kritische toepassingen \autocite{IBMzOSOverview}.

\textcolor{red}{De concrete probleemstelling volgt als:}
\textcolor{red}{
	Het beheer van z/OS-omgevingen binnen grote organisaties zoals banken en overheidsinstellingen vereist sterk gespecialiseerde en praktijkgerichte kennis. 
	Door vergrijzing en een beperkte instroom van nieuwe mainframe-profielen wordt kennisoverdracht vaak impliciet en persoonsafhankelijk georganiseerd, wat leidt tot trage onboarding en verhoogde operationele risico’s. 
	Daarnaast ontbreekt in veel organisaties en onderwijscontexten een eenduidige en gestructureerde afbakening van taken, verantwoordelijkheden en vereiste competenties tussen System Administrator Level~1 en Level~2 binnen z/OS-omgevingen. 
	Het ontbreken van een dergelijk competentieframework bemoeilijkt gerichte opleiding, objectieve evaluatie en doorgroei van startende z/OS System Administrators.}

Bestaande generieke competentiekaders zoals SFIA en e-CF bieden een breed referentiekader voor skills en niveaus, maar blijven vaak te generiek voor de specifieke context van z/OS system administration \autocite{SFIA9,eCF2020}. 
Ook rol en workforce kaders zoals het NIST NICE Framework tonen hoe competenties gestructureerd kunnen worden via rollen, taken en vaardigheidscomponenten, wat relevant is voor de opbouw van een meetbaar framework \autocite{NISTNICE800181r1}. 
Domeinspecifieke kaders bestaan eveneens, zoals IBM-competentie- en apprenticeshipframeworks voor mainframe system administrators (Level~1 en Level~2), maar deze zijn niet altijd uitgewerkt als onderwijs- en onboardinggericht competentieframework met concrete leeractiviteiten en expliciete, meetbare indicatoren \autocite{IBMZSysAdminLevel1,IBMZSysAdminLevel2}.

\textcolor{blue}{De centrale onderzoeksvraag van deze bachelorproef luidt als volgt:}

\textcolor{blue}{\textit{“Welke kerncompetenties zijn noodzakelijk om de rol van System Administrator Level~1 en Level~2 effectief te kunnen invullen binnen een z/OS omgeving?”}}

Om deze onderzoeksvraag te beantwoorden, worden de volgende deelvragen onderzocht.

\textcolor{purple}{\textbf{Deelvragen met betrekking tot het probleemdomein:}}
\textcolor{purple}{
	\begin{itemize}
	\item Welke taken en operationele verantwoordelijkheden (operaties, incident/change, monitoring, batch, security, automatisering) worden binnen z/OS-omgevingen verwacht van System Administrators in organisaties zoals banken en financiële diensten?
	\item Welke verantwoordelijkheden en beslissingsruimte verschillen tussen Level~1 en Level~2 (bv.\ triage vs.\ root-cause, standaard changes vs.\ complexe changes)?
    \item Welke knelpunten en risico's ervaren organisaties wanneer rol- en competentieafbakening ontbreekt (bv.\ escalaties, fouten in changes, compliance-risico)?
    \item Welke leer en onboardingnoden worden het vaakst genoemd bij startende z/OS System Administrators?
\end{itemize}
}

\textcolor{orange}{\textbf{Deelvragen met betrekking tot het oplossingsdomein:}}
\textcolor{orange}{
		\begin{itemize}
		\item Hoe kunnen competenties voor System Administrators binnen z/OS-omgevingen gestructureerd worden in een overzichtelijk framework?
		\item Welke bestaande competentie- en rolframeworks zijn bruikbaar als basis (bv.\ SFIA, e-CF, NICE, IBM Z frameworks)?
		\item Op welke manier kan een onderscheid tussen Level~1 en Level~2 worden uitgewerkt binnen dit framework?
		\item Welke minimale requirements moet een PoC-website bevatten om dit framework bruikbaar te ontsluiten?
		\item Hoe kunnen framework en PoC gevalideerd worden met stakeholders (industrie + onderwijs) op bruikbaarheid, volledigheid en duidelijk onderscheid tussen levels?
	\end{itemize}
}
\textcolor{green}{
	De doelgroep van deze bachelorproef bestaat uit studenten met een focus op mainframe technologie aan HOGENT of andere onderwijsinstellingen met mainframe-gerelateerde opleidingsonderdelen, evenals uit junior of startende System Administrators die actief zijn binnen organisaties waar gebruik wordt gemaakt van IBM mainframes en/of z/OS.
}

In deze bachelorproef wordt het competentieframework inhoudelijk afgebakend en ge\"orienteerd op de IBM Level~1 en Level~2 competency- en apprenticeshipframeworks, zodat de scope aansluit bij de industrieel gehanteerde competenties \autocite{IBMZSysAdminLevel1,IBMZSysAdminLevel2}.

Het concrete eindresultaat van deze bachelorproef is een proof of concept (PoC) in de vorm van een website waarin het competentieframework op een overzichtelijke manier wordt voorgesteld. Met de feedback van een vak-expert (co-promotor) zal de PoC beoordeeld  en verbeterd worden. Het onderzoek wordt als succesvol beschouwd wanneer het framework voldoende duidelijk, herkenbaar en bruikbaar is voor professionals/studenten die actief zijn binnen een z/OS context.


%---------- Stand van zaken ---------------------------------------------------

\section{Literatuurstudie}%
\label{sec:literatuurstudie}

Mainframes blijven een essentieel onderdeel van de IT infrastructuur van grote organisaties, met name in sectoren waar betrouwbaarheid, schaalbaarheid en beveiliging cruciaal zijn zoals banken, overheidsinstellingen en verzekeringsmaatschappijen.
Binnen deze omgevingen speelt het besturingssysteem z/OS een centrale rol. z/OS is specifiek ontworpen om grote hoeveelheden transacties gelijktijdig te verwerken en biedt uitgebreide ondersteuning voor beveiliging, workloadmanagement en systeemautomatisering \autocite{IBMZArchitecture}. Deze eigenschappen maken het beheer van z/OS omgevingen complex en kennisintensief.

\subsection{De rol van de System Administrator (Level 1 en Level 2)}

Binnen IBM mainframe-omgevingen wordt de rol van de z/OS System Administrator onderscheiden in ervaringsniveaus, met name Level~1 (instapniveau) en Level~2 (gevorderd niveau). Elk niveau wordt gekenmerkt door een toenemende mate van autonomie, taakcomplexiteit en impact op de productieomgeving \autocite{IBMZSysAdminLevel1,IBMZSysAdminLevel2}.

Een \textit{System Administrator Level~1} richt zich voornamelijk op operationele en ondersteunende taken binnen de z/OS-omgeving. Volgens het IBM Z System Mainframe Administrator Level~1 Competency Framework omvat dit onder andere het uitvoeren van standaardoperaties, het monitoren van systemen, het ondersteunen bij incidenten, het uitvoeren van vooraf gedefinieerde changes, en het assisteren bij batchverwerking, securityconfiguratie en automatisering \autocite{IBMZSysAdminLevel1}. Deze rol vereist een solide basiskennis van z/OS-concepten en procedures, maar werkt doorgaans onder begeleiding en volgens vastgelegde processen.

De rol van \textit{System Administrator Level~2} bouwt voort op deze operationele basis en omvat een breder en complexer takenpakket. Het IBM zSystems Administrator Level~2 Apprenticeship Framework beschrijft dat Level~2-profielen verantwoordelijk zijn voor het zelfstandig analyseren en oplossen van complexere incidenten, het uitvoeren en beoordelen van niet-standaard changes, het optimaliseren van systeemprestaties en het actief bijdragen aan security-, continuïteits- en recoverystrategieën \autocite{IBMZSysAdminLevel2}. Daarnaast wordt van Level~2 System Administrators verwacht dat zij een coachende rol opnemen ten aanzien van junior profielen en bijdragen aan documentatie en kennisoverdracht binnen het team.

Het expliciete onderscheid tussen beide niveaus is essentieel voor gerichte opleiding, onboarding en doorgroei binnen z/OS-teams en vormt een belangrijke basis voor het competentieframework dat in deze bachelorproef wordt ontwikkeld.


\subsection{Competenties en competentieframeworks in IT}

Competenties worden in de literatuur omschreven als een combinatie van kennis, vaardigheden en professionele attitudes die nodig zijn om een functie effectief uit te voeren. In IT-organisaties worden competentieframeworks gebruikt om deze elementen te structureren en te koppelen aan rollen, verantwoordelijkheden en opleidingsdoelstellingen \autocite{ITILFoundation}.

Generieke frameworks zoals SFIA en het European e-Competence Framework (e-CF) bieden een breed en niveaugebaseerd overzicht van IT-competenties, maar blijven bewust technologie-onafhankelijk \autocite{SFIA9,eCF2020}. Voor gespecialiseerde en bedrijfskritische omgevingen zoals z/OS volstaat deze generieke benadering echter niet. De complexiteit en kriticiteit van mainframeomgevingen vereisen een domeinspecifiek competentieframework dat expliciet rekening houdt met taken, verantwoordelijkheden en ervaringsniveaus \autocite{IBMRedbookzOS,IBMABCsZOSV1}.


\subsection{Kennisoverdracht en ervaringsniveaus binnen IT teams}

Het expliciet definiëren van ervaringsniveaus, zoals System Administrator Level~1 en Level~2, kan bijdragen tot een duidelijker onderscheid in verantwoordelijkheden, verwachtingen en opleidingsbehoeften. Dergelijke niveaudifferentiatie ondersteunt niet alleen kennisoverdracht, maar faciliteert ook gerichte coaching en doorgroei binnen teams.

De bestaande literatuur en documentatie bieden echter weinig concrete richtlijnen voor het structureel invullen van deze ervaringsniveaus binnen z/OS-omgevingen. Hierdoor blijven rolverdelingen vaak impliciet en afhankelijk van organisatie-specifieke afspraken, wat de nood aan een gestructureerd competentiekader verder onderstreept \autocite{IBMzOSOverview}.

% Voor literatuurverwijzingen zijn er twee belangrijke commando's:
% \autocite{KEY} => (Auteur, jaartal) Gebruik dit als de naam van de auteur
%   geen onderdeel is van de zin.
% \textcite{KEY} => Auteur (jaartal)  Gebruik dit als de auteursnaam wel een
%   functie heeft in de zin (bv. ``Uit onderzoek door Doll & Hill (1954) bleek
%   ...'')

\subsection{Synthese en positionering van het onderzoek}

Uit de geraadpleegde literatuur blijkt dat main\- frame en z/OS-omgevingen gekenmerkt worden door een hoge mate van complexiteit en een sterke afhankelijkheid van gespecialiseerde kennis \autocite{IBMABCsZOSV1}. De rol van de System Administrator binnen deze context is essentieel voor het waarborgen van stabiliteit, veiligheid en continuïteit van bedrijfskritische systemen.

Hoewel algemene IT kaders en documentatie het belang van competenties, kennisoverdracht en rolafbakening erkennen, ontbreekt een expliciet en gestructureerd competentiemodel dat rekening houdt met verschillende ervaringsniveaus binnen z/OS-omgevingen. Deze leemte bemoeilijkt zowel opleiding als doorgroei van System Administrators.

Deze bachelorproef positioneert zich binnen deze vastgestelde lacune door te focussen op het expliciteren en structureren van kerncompetenties voor System Administrator Level~1 en Level~2 binnen z/OS-omgevingen. Het beoogde resultaat is een toepasbaar competentieframework dat aansluit bij de noden van zowel organisaties als IT professionals.


%---------- Methodologie ------------------------------------------------------
\section{Methodologie}
\label{sec:methodologie}

Deze bachelorproef hanteert een ontwerpend onderzoeksopzet (design science research) met als doel het ontwikkelen en valideren van een competentieframework voor z/OS System Administrators Level~1 en Level~2. Het onderzoek combineert literatuuronderzoek, praktijkgericht kwalitatief onderzoek en iteratief ontwerp, en wordt ondersteund door feedback van een vakexpert (co-promotor), praktijkstakeholders (KBC/IBS Markets) en domeinspecifieke documentatie (IBM Redbooks). 

De methodologie is bestaat uit vijf fasen.


\subsection{Scope (afbakening t.o.v.\ IBM Level~1/Level~2)}
De inhoudelijke scope van dit onderzoek wordt afgebakend op basis van de IBM competency- en apprenticeshipframeworks voor \textit{Mainframe System Administrator} Level~1 en Level~2 \autocite{IBMZSysAdminLevel1,IBMZSysAdminLevel2}. 
De competenties worden niet \'e\'en-op-\'e\'en overgenomen, maar geclusterd tot een onderwijs- en onboardinggericht framework.

\textbf{In scope} vallen de IBM-competentieclusters rond:
\begin{itemize}
	\item(i) professionele skills (teamwork, communicatie, feedback, leren), 
	\item(ii) change/process/documentatie,
    \item(iii) z/OS systeemcomponenten en interacties (ISPF, TSO commandos, SDSF),
	\item(iv) security concepten (RACF),
	\item(v) batch concepten, JCL en basis scripting/REXX,
	\item(vi) probleemherkenning en operationele troubleshooting op hoofdlijnen
\end{itemize}

\textbf{Out of scope} vallen diepgaande specialistrollen en subsystemexpertise (bv.\ volledige CICS, IMS, Db2 roluitvoering of advanced performance engineering). 
Deze onderwerpen worden enkel vermeld als context omdat ze in de IBM-lijsten voorkomen, maar ze worden niet tot leerdoelen uitgewerkt binnen het eindframework.

\subsection{Fase 1: Literatuurstudie}

In de eerste fase wordt een systematische literatuurstudie uitgevoerd met focus op vier categorieën:
\begin{itemize}
	\item \textbf{Competentieframeworks en pedagogische literatuur}: SFIA, European e-Competence Framework (e-CF), NIST NICE Framework, IBM Z System Administrator frameworks (Level~1 en Level~2) en competentiegericht leren en constructive alignment ter ondersteuning van onderwijs en onboardingtoepassingen.
	\item \textbf{Mainframe en z/OS-literatuur}: IBM Redbooks en officiële z/OS-documentatie met betrekking tot systeembeheer, workload management en operationele verantwoordelijkheden.
	\item \textbf{Mainframe career- en competency studies}: publicaties over skills gaps, vergrijzing en instroomproblematiek binnen mainframe omgevingen.
\end{itemize}

Deze fase resulteert in een eerste set competentiedomeinen en niveau-indicatoren die als input dienen voor het verdere frameworkontwerp.

\subsection{Fase 2: Praktijkonderzoek (interviews)}

Om het probleemdomein te concretiseren en te valideren wordt kwalitatief onderzoek uitgevoerd via semigestructureerde interviews met praktijkstakeholders.
Er wordt minstens één interview gepland met medewerkers actief in een z/OS-context binnen de financiele sector, meer bepaald binnen het \textit{International Banking Systems (IBS)}-team van KBC (Markets).
Deze stakeholders vertegenwoordigen zowel de industriële als de opleidingscontext waarin het competentieframework zal worden toegepast.

De interviews richten zich op:
\begin{itemize}
	\item het dagelijkse takenpakket van z/OS System Administrators;
	\item ervaren knelpunten bij onboarding en kennisoverdracht;
	\item vereiste competenties vanuit praktijkperspectief.
\end{itemize}

De interviewdata wordt geanalyseerd via thematische analyse, waarbij terugkerende patronen en competenties worden geïdentificeerd en vergeleken met de bevindingen uit de literatuur.

\subsection{Fase 3: Requirementsanalyse en ontwerp van het competentieframework (Bachelorproef)}

In deze fase worden de vereisten voor het competentieframework expliciet afgeleid uit de literatuurstudie en de interviewresultaten. Deze requirements bepalen onder meer welke competentiedomeinen worden opgenomen, hoe het onderscheid tussen Level~1 en Level~2 wordt uitgewerkt en welke inhoud minimaal ontsloten moet worden via de proof of concept website.

\subsection{Fase 4: Proof of Concept (PoC)}

Als demonstratie van de toepasbaarheid wordt een proof of concept ontwikkeld in de vorm van een website.
Deze website presenteert het competentieframework op een gestructureerde manier en bevat per competentiedomein:
\begin{itemize}
	\item een theoretische toelichting;
	\item een duidelijke Level~1/Level~2-afbakening;
	\item voorbeeldleeractiviteiten en oefenvormen (meerkeuze vragen, open vragen, quiz).
\end{itemize}

\subsection{Fase 5: Validatie}

Het ontwikkelde competentieframework en de PoC worden gevalideerd via:
\begin{itemize}
	\item feedback van de co-promotor (expert in het vakdomein).
	\item feedback van praktijkstakeholders (o.a.\ interviewdeelnemers).
	\item vergelijking met bestaande IBM-frameworks voor Level~1 en Level~2 (\autocite{IBMZSysAdminLevel1} en \autocite{IBMZSysAdminLevel2})
	\item evaluatie van duidelijkheid en bruikbaarheid voor de doelgroep.
\end{itemize}

Deze validatie waarborgt dat het framework zowel theoretisch onderbouwd als praktisch relevant is.

\subsection{Gebruikte tools en technologieën}

Voor de uitvoering van dit onderzoek worden volgende tools en technologieën gebruikt:
\begin{itemize}
	\item \textbf{Literatuurbeheer}: JabRef in combinatie met \LaTeX{} en Bib\LaTeX{};
	\item \textbf{Kwalitatieve analyse}: semigestructureerde interviews en thematische analyse;
	\item \textbf{Proof of Concept}: HTML, CSS en JavaScript;
	\item \textbf{Versiebeheer en publicatie}: GitHub en GitHub Pages.
\end{itemize}

\subsection{Planning en deliverables}

Deze bachelorproef wordt uitgevoerd binnen het academiejaar 2025--2026 met volgende planning:
\begin{itemize}
	\item \textbf{Februari 2026}: literatuurstudie (z/OS, bestaande competentieframeworks) en opstellen eerste domeinstructuur;
	\item \textbf{Maart 2026}: praktijkinput via (minstens) één semigestructureerd interview met stakeholders uit een z/OS-context (bv.\ KBC/IBS Markets) en verwerking van de resultaten en meer inhoudelijke informatie over Sys Admin verzamelen met de hulp van IBM Redbooks;
	\item \textbf{April 2026}: uitwerking van het competentieframework (inhoudelijk) binnen de bachelorproef met niveau's Level 1 en Level 2.
	\item \textbf{Maart--Mei 2026 (parallel en continu)}: ontwikkeling van een proof of concept website die het framework ontsluit;
	\item \textbf{Februari--Mei 2026 (parallel en continu)}: validatie (stakeholder feedback) en finale rapportering (in Mei).
\end{itemize}

\textbf{Deliverables:}
\begin{itemize}
	\item \textbf{Competentieframework (Bachelorproef)} voor z/OS System Administrators Level~1 en Level~2 (gestructureerd per domein en voorzien van Level-onderscheid);
	\item \textbf{Proof of Concept website}, in de vorm van een informatieve en educatieve webpagina, waarin het competentieframework op een overzichtelijke manier wordt gepresenteerd, met ondersteunende toelichting en voorbeeldoefeningen die het onderscheid tussen Level~1 en Level~2 illustreren.
\end{itemize}


%---------- Verwachte resultaten ----------------------------------------------
\section{Verwachte resultaten en conclusie}%
\label{sec:verwachte_resultaten}

Het verwachte resultaat van deze bachelorproef is de ontwikkeling van een gestructureerd en toepasbaar competentieframework voor z/OS System Administrators Level~1 en Level~2. Dit framework brengt de kerncompetenties in kaart die nodig zijn om deze rollen effectief uit te voeren, met aandacht voor technische kennis, operationele vaardigheden en professionele verantwoordelijkheden, en maakt het onderscheid tussen beide niveaus expliciet.

Daarnaast wordt een proof of concept ontwikkeld in de vorm van een website die het competentieframework overzichtelijk ontsluit en illustreert met ondersteunende toelichting en voorbeeldleeractiviteiten. De meerwaarde van dit onderzoek ligt in het expliciteren en structureren van kennis binnen een gespecialiseerd domein, zodat onboarding en opleiding gerichter kunnen verlopen en kennisoverdracht minder persoonsafhankelijk wordt. 
